\section{Background and related work}
\label{sec:background}
\subsection{Learning from irregular clinical observations}
A hospital record gives the time of a measurement more readily than the duration of the state it captures. GRU-D uses both missingness and the time elapsed since the last measurement \citep{che2018recurrent}. STraTS encodes the variable, value, and time of each observation and learns how observations relate to one another \citep{tipirneni2022strats}. Both can therefore use measurements taken at irregular intervals.

Knowledge-based temporal abstraction (KBTA) defines some of these relationships before model fitting \citep{shahar1997framework}. The Temporal Mediator applies the rules in a clinical knowledge base to time-stamped source concepts such as glucose measurements \citep{mediator-github}. It derives states, trends, events, and contexts, which may hold at a point in time or over an interval. The same glucose readings can contribute to two overlapping intervals, one indicating that glucose is high and the other that it is falling. Earlier methods identified recurring relationships between clinical intervals and used their frequencies and durations to predict outcomes \citep{moskovitch2015karmalego,sheetrit2019tpf}. INTERVenE represents the beginning and end of each interval separately, allowing the encoder to learn from duration and overlap while retaining the clinical identity of each interval \citep{oded2026intervene}.

\subsection{Representing guideline adherence over time}
The Kineret clinical study used the Temporal Mediator to assess fourteen components of inpatient diabetes care \citep{oded2027kineret-jbi,hochberg2023inpatient}. The Mediator then evaluates quality-assessment temporal patterns (QATPs) over the resulting clinical history. Each quality-assessment (QA) pattern specifies when a recommendation applies and what action is expected. For insulin after hyperglycemia, the assessment asks whether insulin was given during the period specified by the rule. For routine monitoring, the score depends on how often glucose was checked within a specified interval. These rules allow partial credit and can capture repeated assessments during a stay.

Two patients may have the same mean compliance score even if a missed or delayed action occurs at different points in their glucose trajectories. Kineret records both the scores and the intervals to which they refer. A model can therefore relate a particular assessment to the patient's course as well as use a summary of adherence up to the time of prediction.
